%!TEX root = ./main.tex

\section[Act I: Space]{Act I: A Cell Tower in Space}

% SECTION TIMING: 7:00--25:00 (8 frames). Paper A (SN2, frames 1-5) runs 7:00-17:00;
% Paper B (StarCDN, frames 6-8) runs 17:00-25:00. Both \watch clips are 60-90 s; if
% discussion runs long, cut the StarCDN clip, not the student-question frame.
% LAYOUT RULE: every bullet must fit on one line. Shorten the text, do not let it wrap.

% ---------------------------------------------------------------------------
% PAPER A: SN2 -- Direct-to-Cell Satellite Network without Satellite Navigation
% ---------------------------------------------------------------------------

% Teaching note (2 min): ask who has seen a Starlink dish. Then land the twist:
% this paper is about the version with NO dish -- the phone already in your pocket.
% The talk clip moved to the "let the network navigate itself" frame -- it lands
% better after the GPS twist than before it.
% SOURCE (commercial D2C service is real; launched July 23, 2025; texting on
% ordinary recent phones, data added Oct 2025):
% https://broadbandbreakfast.com/t-mobile-and-starlink-satellite-service-to-officially-launch-july-23/
% https://www.t-mobile.com/coverage/satellite-phone-service
\begin{frame}{The tower left the ground}
  \large\textbf{Classic Starlink needs a dish on your roof.}\\[0.3em]
  What if your ordinary phone talks straight to the satellite?

  \vspace{0.4em}
  \small
  Not hypothetical: T-Mobile + Starlink have sold satellite texting\\
  on unmodified recent phones since July 2025.

  % SOURCE (LEO orbital speed ~7.5-7.8 km/s is standard orbital mechanics):
  % https://en.wikipedia.org/wiki/Low_Earth_orbit
  \brokenassumption{The base station stands on the ground}
    {a cell tower moving at about 7.5 km/s in LEO}

  \vspace{0.4em}
  \papercard{Direct-to-Cell Satellite Network without Satellite Navigation}
    {Wei Liu, Yuanjie Li, Jingyi Lan, Hewu Li, et al. --- Tsinghua University}
    {ACM SIGCOMM 2025}
    {https://doi.org/10.1145/3718958.3750522}

  \glossline{LEO = Low Earth Orbit \quad D2C = direct-to-cell}
\end{frame}

% Teaching note (2.5-3 min, HARD CAP 3 -- the section must fit 18:00): after
% showing the diagram, run the \ask below and let
% students shout research questions. They reliably produce: handoff every few
% minutes, Doppler, "which country's satellite serves me", authentication,
% billing, location. Write them on the board -- Act I and Act III both pay
% several of them off. Do NOT answer them yet.
\begin{frame}{Same phone. Very different tower.}
  \centering
  \begin{tikzpicture}[scale=0.92,transform shape]
    % Left: today
    \node[netnode,minimum width=1.5cm] (ph1) at (0,0) {Phone};
    \node[netnode,minimum width=1.6cm] (tw) at (3.1,0) {Tower};
    \draw[flow] (ph1) -- (tw);
    \node[font=\scriptsize\itshape,align=center] at (1.55,-0.75)
      {bolted to the ground\\yours for years};
    % Right: D2C
    \node[netnode,minimum width=1.5cm] (ph2) at (6.4,0) {Phone};
    \node[oranode,minimum width=1.7cm] (sat) at (9.5,1.1) {Satellite};
    \draw[flow] (ph2) -- (sat);
    \draw[flow,draw=ranpurple,dashed] (sat) -- ++(1.35,0)
      node[right,font=\scriptsize,text=ranpurple] {7.5 km/s};
    \node[font=\scriptsize\itshape,align=center] at (8.0,-0.75)
      {overhead for minutes\\then gone};
  \end{tikzpicture}

  \vspace{0.7em}
  \ask{Your base station now moves at 7.5 km/s.\\
       Sixty seconds: shout out everything that breaks.}

  \vspace{0.3em}
  {\small\textcolor{softgray}{Keep your list. This lecture pays several of
   these off.}\par}
\end{frame}

% Teaching note (2 min): the reveal. Students expect Doppler or handoff to be
% the paper's problem; the actual problem is a dependency nobody in the room
% will have named. Pause before the last line.
% SOURCE (D2C relies on GNSS for radio access, authentication, authorization;
% GNSS defects propagate into the D2C network):
% https://doi.org/10.1145/3718958.3750522
\begin{frame}{The twist: it secretly runs on GPS}
  \large\textbf{Today's D2C networks lean on GNSS \expand{GPS and friends}.}

  \vspace{0.5em}
  \small
  \begin{itemize}\itemsep5pt
    \item \textbf{Where you are} decides beam, country, and bill.
    \item \textbf{What time it is} decides when your radio may talk
      \expand{uplink slots run on a shared clock}.
    \item Both answers come from a \emph{different} satellite system.
  \end{itemize}

  \vspace{0.3em}
  A ground tower knows where it is. A moving one has to keep asking.

  \vspace{0.3em}
  But GNSS can be \textbf{jammed}, \textbf{spoofed}, or simply
  \textbf{unavailable}.

  % SOURCE (observed failures: intermittent connectivity, over/under-billing,
  % unauthorized service, service denials even when D2C sats are reachable):
  % https://doi.org/10.1145/3718958.3750522
  \qa{Satellite bars on your phone, but no GPS fix. What happens?}
     {Field tests: denials, dropped access, wrong bills ---
      \textbf{even with the D2C satellite overhead}. One satellite
      system's correctness hangs on another one.}

  \glossline{GNSS = Global Navigation Satellite System (GPS, Galileo, BeiDou, ...)\\
    beam = the patch of ground the satellite aims at \quad
    uplink = the phone-to-satellite direction}
\end{frame}

% Teaching note (2 min): keep this at the idea level. Do NOT derive positioning
% from Doppler or delay -- one sentence of "the signals you already receive
% carry enough geometry" is the ceiling for this lecture.
% SOURCE (fate-sharing principle; reuse D2C satellites as the navigation
% source; trade navigation accuracy for network availability):
% https://doi.org/10.1145/3718958.3750522
\begin{frame}{The move: let the network navigate itself}
  \large\textbf{SN2's idea: stop borrowing. Self-navigate.}

  \vspace{0.5em}
  \small
  \begin{itemize}\itemsep5pt
    \item The D2C satellites themselves become the position source.
    \item If you can hear the satellite, you can locate by it.
    \item They live or die together --- the authors call it \textbf{fate-sharing}.
  \end{itemize}

  \qa{Doesn't the network need GPS-grade precision?}
     {No. Beam, country, and bill need \textbf{good-enough} position,
      not meter-level.}

  \watch{https://www.youtube.com/watch?v=G3KDDdRZYBE}{SIGCOMM'25 talk --- play 60--90 s}
\end{frame}

% SOURCE (both ranges below; evaluated with commodity phones with 3GPP NTN
% support, against legacy GNSS-dependent solutions):
% https://doi.org/10.1145/3718958.3750522
% Teaching note (90 s): say aloud what the range endpoints correspond to --
% 4.4x and 23.5x are the extremes across the scenarios the paper tests (mild
% vs severe GNSS unavailability). Verify the exact scenario mapping against
% the paper's evaluation section before class; do not improvise it.
% Teaching note: also say aloud what "network availability" measures here --
% per the paper, the share of access attempts that succeed (it is NOT tower
% uptime). VERIFY the exact definition against the paper's evaluation section
% before class.
\begin{frame}{What refusing to need GPS bought}
  \begin{columns}[T,onlytextwidth]
    \column{0.47\textwidth}
      \large\textbf{Network availability}

      \vspace{0.4em}
      {\Huge\textcolor{ranblue}{4.4--23.5\(\times\)}}

      \vspace{0.3em}
      \small higher than legacy GNSS-dependent designs
    \column{0.47\textwidth}
      \large\textbf{Access latency}

      \vspace{0.4em}
      {\Huge\textcolor{coreorange}{1.9--12.3\(\times\)}}

      \vspace{0.3em}
      \small lower, same comparison
  \end{columns}

  \vspace{0.5em}
  {\scriptsize\textcolor{softgray}{Evaluated with commodity phones with 3GPP
   NTN support. Ranges span the scenarios tested.}\par}

  \glossline{NTN = Non-Terrestrial Network, 3GPP's name for satellite cellular}

  \vspace{0.3em}
  \takeaway{The fix was not better GPS.\\
    It was deleting the hidden dependency on GPS.}
\end{frame}

% ---------------------------------------------------------------------------
% PAPER B: StarCDN -- Moving Content Delivery Networks to Space
% ---------------------------------------------------------------------------

% Teaching note (2 min): start from the students' own evening. Everyone has
% streamed video at peak hours; almost nobody has asked why it works.
% One-liner worth saying at the papercard: the incumbent CDN (Akamai)
% co-authored the paper about replacing ground CDNs.
\begin{frame}{Why is Netflix fast at 8 pm?}
  \qa{Millions stream at once. Why doesn't it collapse?}
     {Because the video rarely comes from far away. CDNs
      \expand{Content Delivery Networks} cache popular content in
      servers \textbf{near you} --- your city, sometimes your ISP.}

  \large\textbf{Now route your Internet through space.}

  \vspace{0.1em}
  \small
  Every request climbs to a satellite, drops to a ground station,\\
  finds the cache on Earth --- then climbs all the way back up.

  \brokenassumption{The cache sits on the ground, near the user}
    {putting the CDN on the satellites themselves}

  \vspace{0.1em}
  {\small So well hidden it is not even a box in your diagram.\par}

  \vspace{0.1em}
  \papercard{StarCDN: Moving Content Delivery Networks to Space}
    {W.\,X. Zheng, A. Taneja, M. Masood, A. Sabnis, R.\,K. Sitaraman,
     D. Vasisht --- UIUC, UMass Amherst \& Akamai}
    {ACM SIGCOMM 2025}
    {https://doi.org/10.1145/3718958.3754345}
  \watch{https://www.youtube.com/watch?v=-dX0DbW-9iY}{SIGCOMM'25 talk --- play 60--90 s}
\end{frame}

% Teaching note (3 min): this is the frame students remember. Let them sit with
% the problem before revealing the arrow trick: the satellites drift one way,
% the content is handed the other way, so the CONTENT hovers over Florida even
% though no SATELLITE does.
% SOURCE (LEO orbital period ~90-100 min is standard orbital mechanics):
% https://en.wikipedia.org/wiki/Low_Earth_orbit
\begin{frame}{Your cache is doing 7.5 km/s}
  \small
  A LEO satellite laps the Earth in roughly 95 minutes.\\
  The one over Florida right now is Florida's for only minutes.

  \qa{It spent those minutes caching what Floridians watch. Then what?}
     {StarCDN hands the cached content \textbf{backward} through the
      constellation \expand{the whole fleet} --- against the satellites'
      motion --- to the next satellite arriving overhead.}

  \vspace{0.15em}
  % SOURCE (the constellation's inter-satellite laser links are the transport
  % that carries content between satellites in StarCDN):
  % https://doi.org/10.1145/3718958.3754345
  Satellites carry laser links to their neighbors --- content rides those.

  \vspace{0.15em}
  \centering
  % Teaching note: the picture is deliberately picture-level -- real orbits are
  % inclined planes, not neat rows. Say that aloud if a sharp student objects.
  \begin{tikzpicture}[scale=0.88,transform shape]
    \node[oranode,minimum width=1.4cm] (s1) at (0,0) {Sat A};
    \node[oranode,minimum width=1.4cm] (s2) at (3.0,0) {Sat B};
    \node[oranode,minimum width=1.4cm] (s3) at (6.0,0) {Sat C};
    \node[font=\scriptsize,text=softgray,anchor=south] at (0,0.42) {next up};
    \node[font=\scriptsize,text=softgray,anchor=south] at (3.0,0.42) {serving now};
    \node[font=\scriptsize,text=softgray,anchor=south] at (6.0,0.42) {just left};
    \draw[flow,draw=ranpurple,dashed] (7.1,0.65) -- (8.6,0.65)
      node[right,font=\scriptsize,text=ranpurple] {satellites drift};
    \draw[flow,draw=coreorange,very thick] (s3) -- (s2);
    \draw[flow,draw=coreorange,very thick] (s2) -- (s1);
    \node[font=\scriptsize,text=coreorange] at (3.0,-0.75)
      {content handed backward};
    \node[netnode,minimum width=2.2cm] (fl) at (3.0,-1.9) {Users in Florida};
    \draw[flow,dashed] (fl) -- (s2);
  \end{tikzpicture}

  \vspace{0.3em}
  {\small A useful way to picture it: \emph{the satellite moves on; the
   content stays behind.} (Picture-level, not real geometry.)\par}
\end{frame}

% SOURCE (both numbers below: ~80% reduction in ground-to-satellite bandwidth
% usage; ~2.5x improvement in user-perceived latency):
% https://doi.org/10.1145/3718958.3754345
\begin{frame}{What moving the cache to orbit bought}
  \begin{columns}[T,onlytextwidth]
    \column{0.47\textwidth}
      \large\textbf{Ground--satellite bandwidth}

      \vspace{0.4em}
      {\Huge\textcolor{ranblue}{\(\sim\)80\%}}

      \vspace{0.3em}
      \small less usage on the scarcest link
    \column{0.47\textwidth}
      \large\textbf{User-perceived latency}

      \vspace{0.4em}
      {\Huge\textcolor{coreorange}{\(\sim\)2.5\(\times\)}}

      \vspace{0.3em}
      \small improvement over ground-based CDN baselines
  \end{columns}

  \vspace{0.3em}
  {\scriptsize\textcolor{softgray}{Trace-driven simulation using real-world
   Akamai CDN traces.}\par}

  \vspace{0.3em}
  \takeaway{Act I score: the \textbf{tower} box and the \textbf{CDN} box
    both left the ground --- and the systems got \emph{better}.}

  \vspace{0.4em}
  % The strip below is the section-01 diagram, shrunk, so students can look at
  % the picture the closing question is about.
  \centering
  \begin{tikzpicture}[scale=0.7,transform shape]
    \node[netnode,minimum width=1.5cm]  (ue)   at (0,0)    {UE};
    \node[netnode,minimum width=1.35cm] (ru)   at (2.15,0) {RU};
    \node[netnode,minimum width=1.35cm] (du)   at (4.05,0) {DU};
    \node[netnode,minimum width=1.35cm] (cu)   at (5.95,0) {CU};
    \node[corenode,minimum width=1.6cm] (core) at (8.35,0) {5G Core};
    \draw[flow] (ue) -- (ru);
    \draw[flow] (ru) -- (du);
    \draw[flow] (du) -- (cu);
    \draw[flow] (cu) -- (core);
  \end{tikzpicture}

  \vspace{0.2em}
  \ask{Look at the rest of the architecture diagram.\\
       Which boxes are on the ground only by habit?}
\end{frame}
